Ainsi ces pères mènent leur vie dans l’ombre,\\
Pareils à sept lis, à la parure de neige,\\
De la Vierge très haute fort agréés, à Pierre\\
Ils paraissent briller.\\
\\
Bientôt entraînés par la divine flamme,\\
Ils parcourent les villes, errant en tous lieux !\\
Pour graver, si possible, en tous les cœurs,\\
Les douleurs de leur Mère.\\
\\
C’est ainsi qu’elles peuvent dompter d’aveugles colères,\\
Unir des cœurs farouches, ignorants de la paix,\\
Relever les affligés, ramener les coupables,\\
Leurs pieuses paroles.\\
\\
La Vierge qui accompagnait ses serviteurs\\
Les conduit enfin au rivage céleste,\\
Et de pierres précieuses, pour toute l’éternité,\\
Décore les bienheureux.\\
\\
Et maintenant, le gémissement de la foule priante,\\
Qu’ils l’entendent et voient nos durs labeurs !\\
Et que toujours ils favorisent nos vœux\\
D’un bienveillant regard.\\